\documentclass[11pt]{article}

\usepackage[a4paper,margin=1in]{geometry}
\usepackage[T1]{fontenc}
\usepackage{lmodern}
\usepackage{microtype}
\usepackage{parskip}
\usepackage{xcolor}
\usepackage{booktabs}
\usepackage{tabularx}
\usepackage{enumitem}
\usepackage{titlesec}
\usepackage[hidelinks]{hyperref}

\newcolumntype{Y}{>{\raggedright\arraybackslash}X}

\definecolor{sayc}{HTML}{1F5FBF}
\definecolor{showc}{HTML}{0F766E}
\definecolor{whyc}{HTML}{7C3AED}

\newcommand{\say}{\textcolor{sayc}{\textbf{[SAY]}}\ }
\newcommand{\showcue}{\textcolor{showc}{\textbf{[SHOW]}}\ }
\newcommand{\why}{\textcolor{whyc}{\textbf{[WHY]}}\ }

\titleformat{\section}{\Large\bfseries}{\thesection.}{0.6em}{}
\titleformat{\subsection}{\large\bfseries}{\thesubsection}{0.6em}{}
\setlist[itemize]{leftmargin=1.4em,itemsep=2pt,topsep=2pt}

\title{\textbf{Query Decomposer Arena}\\[2pt]
\large End-to-End Presentation Script}
\author{}
\date{}

\begin{document}
\maketitle
\vspace{-2.5em}

\noindent\textit{Cue legend:}\quad
\say talk track \quad
\showcue what to point at \quad
\why the reasoning, if pressed.\\
\textit{Approx.\ 20--25 min plus Q\&A. Trim Sections~5 and~9 if short on time.}

\hrulefill

\section*{0.\quad One-line pitch (say this first, 20\,s)}
\addcontentsline{toc}{section}{0. One-line pitch}

\say ``We built a benchmarking harness that takes a single complex query, runs it
through nine different query-decomposition methods --- from a dumb rule-based
splitter up to three 2026 SOTA papers --- and scores every method side by side on
how well it breaks the query into sub-questions and a dependency graph. The
contribution isn't a new decomposition algorithm; it's the evaluation framework,
a reference-free scoring rubric, and a normalization layer that makes these very
different methods comparable.''

\section{Problem \& motivation --- the \emph{why} (2 min)}

\say ``When a user asks an LLM a complex question --- something with several
intents, or that needs multi-step reasoning --- you get better, cheaper, faster
answers if you first \emph{decompose} it: break it into smaller sub-questions,
figure out which ones depend on which, then solve that graph.

This matters most in \textbf{edge--cloud LLM systems}. You have a small model
on-device --- cheap, fast, private --- and a large model in the cloud ---
expensive, slow, capable. The decomposition tells you which sub-tasks the small
model can handle and which need the cloud. Get the decomposition wrong and you
either overpay by sending everything to the cloud, or you get wrong answers
because the small model choked on a hard step.

There's a wave of recent work on exactly this --- HybridFlow at ICML~2026,
Route-and-Reason at WWW~2026, Uno-Orchestra --- and they all have a decomposition
component. But they're evaluated on different datasets, with different metrics, so
\textbf{you can't tell which decomposition strategy is actually better}. That's
the gap we're filling.''

\why Heterogeneous methods, no common evaluation $\rightarrow$ our job is to make
them comparable and measure them.

\section{What we built --- the \emph{what} (1 min)}

\say ``A full-stack web app with three parts:''
\begin{enumerate}[leftmargin=1.6em,itemsep=2pt]
  \item \textbf{Catalog} --- a browsable reference of all nine decomposers: how
  each works, what output shape it produces, the source paper.
  \item \textbf{Playground} --- you paste a query, pick which decomposers to run,
  hit Compare. You get every method's sub-questions, an auto-laid-out dependency
  DAG, structural metrics, and a judge score, all side by side.
  \item \textbf{Scoreboard} --- aggregates every comparison you've run into a
  ranked leaderboard with per-criterion breakdowns.
\end{enumerate}

\say ``\textbf{Scope is deliberately narrow: decomposition only.} We produce the
sub-questions and the graph and we score \emph{those}. We do \textbf{not} answer
the sub-questions --- no retrieval, no final-answer generation. That isolates
decomposition quality from every downstream confound.''

\why If you also executed the sub-queries, a bad final answer could be the
retriever's fault, the solver's fault, or the decomposition's fault. By stopping
at the graph, every number we report is about decomposition and nothing else.

\section{The decomposer landscape --- taxonomy (2 min)}

\showcue Catalog page, filter buttons.

\say ``We classify methods on two axes.

\textbf{Axis 1 --- is there an LLM in the loop?}''
\begin{itemize}
  \item \emph{Deterministic}: rule-based, no model. Free, instant, fully
  reproducible.
  \item \emph{LLM-based}: a model produces the decomposition.
\end{itemize}

\say ``\textbf{Axis 2 --- what shape is the output?}''
\begin{itemize}
  \item \emph{Flat}: a bag of independent sub-questions, no dependencies ---
  everything runs in parallel.
  \item \emph{Linear}: a strict chain --- step 2 needs step 1, step 3 needs step 2.
  \item \emph{DAG}: a real dependency graph --- some branches parallel, some
  sequential.
  \item \emph{Adaptive}: the method decides \emph{whether} to decompose at all;
  simple queries stay whole.
\end{itemize}

\say ``These shapes aren't cosmetic. A flat decomposition claims `all these
sub-tasks are independent' --- if that's wrong, you'll parallelize things that
actually needed ordering and get garbage. A linear chain claims `nothing here is
parallelizable' --- if that's wrong, you're leaving latency on the table.
\textbf{The shape is a testable claim about the query.}''

\section{The baselines --- what each is and \emph{why it's there} (2--3 min)}

\say ``You can't say a SOTA method is `good' without reference points. We have a
ladder of baselines, each answering a specific question.''

\begin{tabularx}{\linewidth}{@{}>{\raggedright\arraybackslash}p{2.5cm} Y Y@{}}
\toprule
\textbf{Method} & \textbf{What it is} & \textbf{Why it's in the study} \\
\midrule
Rule-based splitter &
spaCy: split on sentence boundaries and conjunctions (``and'', ``then'', ``;''),
add a sequential edge only when a clause refers back with a pronoun (``it'', ``the
result'') &
The \textbf{zero-cost, zero-LLM floor.} If an expensive LLM method can't beat
regex + a parser, it isn't earning its cost. \\
\addlinespace
Naive prompt --- small model &
One generic instruction: ``break this into atomic sub-questions, note
dependencies'' &
Baseline half of a \textbf{scale ablation.} \\
\addlinespace
Naive prompt --- large model &
\emph{Identical prompt}, large model &
Other half. The gap between these two isolates \textbf{how much decomposition
quality is just model size vs.\ the method/prompt.} This is the key control. \\
\addlinespace
LlamaIndex SubQuestionQueryEngine &
Reproduces the question generator from a production RAG library --- flat JSON list &
\textbf{The practitioner baseline} --- what an engineer actually reaches for today. \\
\addlinespace
Least-to-Most (Zhou et al., 2022) &
Ordered list of easier sub-problems $\rightarrow$ linear chain &
\textbf{Pre-DAG-era classic.} Shows what the field did before dependency graphs. \\
\addlinespace
Self-Ask (Press et al., 2022) &
Chain of follow-up questions, length decided by the model &
Classic, and the one method with \textbf{dynamic decomposition size.} \\
\bottomrule
\end{tabularx}

\why The naive-small vs naive-large pair is the most important design choice.
Without it, a reviewer can say ``your SOTA method just wins because it uses a
bigger model.'' With it, we separate \emph{method contribution} from \emph{scale
contribution}.

\section{The SOTA methods --- how they work (3 min)}

\say ``Then the three frontier methods. We reproduce each paper's
\textbf{decomposition prompt and procedure} for inference only --- no fine-tuning,
no RL checkpoints --- so the comparison is about the \emph{method design}, not who
trained longer.''

\paragraph{HybridFlow (ICML 2026) --- the richest DAG producer.}
\say ``Uses an `Explain--Analyze--Generate' meta-prompt. The planner emits an XML
plan; we deterministically parse it into a graph where each sub-task has a
description, a role label, and a list of prerequisite IDs. Then we \emph{validate}:
is it acyclic? is every node reachable? If not, we re-prompt the model with the
specific error, up to twice. If it still fails, we fall back to a sequential
chain. Capped at 7 sub-tasks, matching the paper. Role `Generate' routes to the
large model, the rest to small.''

\paragraph{Route-and-Reason / R2-Reasoner (WWW 2026) --- linear + routing.}
\say ``Its router has two parts: a Task Decomposer that produces a chain of 3--8
steps, and a Subtask Allocator that assigns each step to a model. We reproduce the
decomposer prompt verbatim from the appendix. For the allocator, which in the
paper is RL-trained, we use a lightweight difficulty heuristic as a documented
stand-in --- steps mentioning `prove', `compare', `optimize' go large, the rest
go small.''

\paragraph{Uno-Orchestra --- parsimonious / adaptive.}
\say ``One LLM call that jointly decides \emph{whether} to decompose and how to
route. Simple queries collapse to a single direct-answer node at zero
orchestration cost; only genuinely compositional queries pay for a multi-node
graph. It's the method that can correctly answer: `don't decompose this.'\,''

\showcue Run a simple query (``What is the capital of Australia?'') and point at
Uno-Orchestra returning one node while the others over-split.

\section{Making them comparable --- normalization \& guarantees (1.5 min)}

\say ``These nine methods produce wildly different raw output --- XML, JSON,
numbered lists, prose. We normalize everything to one schema: \textbf{a list of
sub-queries with IDs, plus a list of directed edges} where an edge
A~$\rightarrow$~B means B depends on A.

Then two guarantees are enforced centrally, so every method gets them for free:''
\begin{enumerate}[leftmargin=1.6em,itemsep=2pt]
  \item \textbf{Always acyclic.} We drop edges that point to non-existent nodes,
  then break any cycles with a minimal edge removal, and flag it. So `is it a
  valid DAG' is never in question --- a method can't cheat by emitting a
  contradictory graph.
  \item \textbf{No invention.} Every LLM method's system prompt carries a shared
  rule: \emph{do not introduce any entity, name, number, or constraint that isn't
  in the query.} This is enforced at the prompt level AND checked independently by
  the judge.
\end{enumerate}

\why Normalization is a real contribution --- it's what lets you put HybridFlow's
XML DAG next to Self-Ask's prose chain on the same axis.

\section{Scoring --- the \emph{how} of measurement (3--4 min)}

\say ``Two kinds of measurement: structural metrics we \emph{compute}, and quality
criteria we \emph{judge}.''

\subsection{Structural metrics --- deterministic, from graph theory}

\say ``These are objective and free. We compute them with a graph library:''
\begin{itemize}
  \item \textbf{Sub-query count} --- how granular is the decomposition.
  \item \textbf{Depth} --- the longest dependency chain. This is the
  \textbf{critical path}: it lower-bounds end-to-end latency, because you can't
  parallelize a chain.
  \item \textbf{Parallelism / width} --- the largest set of mutually independent
  sub-queries (formally, the maximum antichain). This upper-bounds your speedup
  from parallel execution.
  \item \textbf{Roots / leaves} --- entry and exit points of the graph.
\end{itemize}

\say ``We do \textbf{not} ask an LLM to count these. If a computer can measure it
exactly, an LLM shouldn't estimate it.''

\why Depth and width are the metrics that turn a decomposition into a
latency/cost prediction. A method that produces depth~10 is claiming a 10-step
serial pipeline; a method with width~3 is claiming 3$\times$ parallelism is
available.

\subsection{LLM-as-judge --- reference-free, 5 criteria}

\say ``For quality we use an LLM as a judge. Crucially, it's \textbf{reference-free}
--- we score the decomposition against the \emph{original query itself}, not
against a gold-standard decomposition.''

\say ``\textit{Why reference-free:} there is no gold-standard dataset of `correct'
decompositions for arbitrary queries. Building one is expensive, and worse, it's
subjective --- reasonable people decompose the same query differently. So instead
of asking `does this match the one true answer', we ask `is this decomposition
internally sound with respect to the query'. Five criteria, each scored 1 to 5:''

\begin{tabularx}{\linewidth}{@{}>{\raggedright\arraybackslash}p{3.1cm} Y Y@{}}
\toprule
\textbf{Criterion} & \textbf{Question it asks} & \textbf{Why it matters} \\
\midrule
Coverage &
Is every intent and constraint in the query captured by some sub-query? &
Recall. Missing an intent means the final answer will be incomplete. \\
\addlinespace
Minimality / non-redundancy &
Are sub-queries non-overlapping and not over-split? &
Redundant sub-queries waste model calls; over-splitting inflates latency and cost. \\
\addlinespace
Faithfulness &
Did it avoid inventing entities, numbers, or constraints not in the query? &
Hallucinated constraints steer the whole pipeline wrong. Score~1 if anything is
fabricated. \\
\addlinespace
Standalone answerability &
Can each sub-query be answered given only the answers to its declared parents? &
If a sub-query secretly needs context it doesn't list, the dependency graph is a
lie. \\
\addlinespace
Dependency correctness &
Are the parallel-vs-sequential edges right --- none missing, none spurious? &
\textbf{This is what makes the latency and cost numbers trustworthy.} Wrong edges
= wrong performance claims. \\
\bottomrule
\end{tabularx}

\say ``Overall score is the mean of the five. Per run, we judge each of the nine
decompositions once. The Scoreboard then averages across every run in the session
to give a ranked leaderboard plus the per-criterion breakdown and the average
structural stats.''

\showcue Scoreboard: point at the ranking, then at a per-criterion column where a
method is weak (e.g.\ Least-to-Most low on minimality because it over-splits).

\section{Architecture --- the \emph{how} of the system (1.5 min)}

\say
\begin{itemize}
  \item \textbf{Backend}: FastAPI, completely stateless --- no database. One
  endpoint runs all selected decomposers concurrently, isolates failures per
  method, then judges each result.
  \item \textbf{LLMs}: Hugging Face Inference Providers. We pass a \emph{pool} of
  API tokens; the backend round-robins across them and rotates on rate-limit
  errors, so a long comparison doesn't stall. Two model tiers, `small' and
  `large', both configurable --- nothing is hard-coded.
  \item \textbf{Judge model}: defaults to the large model. We're upfront that this
  creates a mild self-bias toward large-tier methods; it's flagged in the UI and
  the judge model is swappable.
  \item \textbf{Frontend}: React + Vite + TypeScript. The DAG visualization is
  React Flow with automatic layout; the judge scorecards are radar charts.
  \item \textbf{Tests}: 20 unit tests, all with the LLM stubbed --- parser
  robustness, graph-metric correctness, and the HybridFlow repair-and-fallback
  path.
\end{itemize}

\section{Live demo (2 min) --- exact steps}

\say ``Let me run one query.'' \showcue Paste:

\begin{quote}
\itshape
``I'm designing an edge--cloud LLM system where a small on-device model and a
large cloud model work together. Compare how HybridFlow, Route-and-Reason, and
Uno-Orchestra each decompose a user query into subtasks and decide which subtasks
go to the small model versus the large one. Then, assuming a budget of at most two
cloud calls per query and a latency target under two seconds, tell me which
approach is the best fit and the main trade-off.''
\end{quote}

\say ``This query is deliberately compositional --- three method descriptions that
can be worked out in parallel, then a comparison that depends on all three, then a
trade-off that depends on the comparison.''

\showcue In order:
\begin{enumerate}[leftmargin=1.6em,itemsep=2pt]
  \item \textbf{HybridFlow / Uno-Orchestra panels} --- point at the DAG: width~3
  at the top (the three methods, parallel), narrowing to a chain (compare
  $\rightarrow$ pick $\rightarrow$ trade-off). ``This is the correct shape.''
  \item \textbf{LlamaIndex / rule-based panels} --- ``Flat. It lost the ordering
  entirely --- it would try to answer `which is best' before describing the
  methods.''
  \item \textbf{naive-small vs naive-large} --- ``Same prompt. The small model
  over-splits into a dozen shallow questions; the large model produces a clean
  graph. That gap is the scale effect.''
  \item \textbf{Judge radars} --- point at dependency-correctness: high for the
  DAG methods, low for the flat ones.
  \item \textbf{Scoreboard} --- ``Ranked. The DAG-native methods lead on
  dependency correctness; the flat methods lead on latency because they're one
  model call.''
\end{enumerate}

\section{What the comparison reveals --- findings (1.5 min)}

\say
\begin{itemize}
  \item \textbf{Model scale matters, but so does method.} naive-large already
  beats the classic methods on our rubric --- but the DAG-native methods
  (HybridFlow, Uno) beat naive-large on dependency correctness at the \emph{same}
  model size.
  \item \textbf{Flat methods systematically lose dependency correctness} on
  compositional queries --- they can't represent ordering.
  \item \textbf{The classic linear methods (Least-to-Most, Self-Ask)
  over-decompose} --- dinged on minimality because they force a chain even when
  branches are independent.
  \item \textbf{Adaptive decomposition is underrated} --- Uno-Orchestra is the
  only method that correctly \emph{declines} to decompose atomic queries; the
  rest always split, wasting calls.
  \item \textbf{The rule-based floor is not zero} --- for simple ``A and B''
  queries it's competitive with the LLMs at zero cost, which tells you when
  \emph{not} to spend a model call on decomposition.
\end{itemize}

\section{Limitations --- be honest (1.5 min)}

\say
\begin{itemize}
  \item \textbf{Judge self-bias.} Judge = large model, so large-tier methods may
  get a small unearned edge. Flagged; mitigation would be an ensemble of judges or
  a different judge family.
  \item \textbf{Reference-free scoring is an approximation.} We measure internal
  soundness, not correctness against a ground truth, because no ground truth
  exists for arbitrary queries.
  \item \textbf{We reproduce prompts, not trained systems.} R2-Reasoner's real
  allocator is RL-trained; HybridFlow's planner is fine-tuned. Our versions use
  base instruct models, so absolute numbers understate the papers --- but the
  \emph{relative method comparison} holds.
  \item \textbf{Decomposition only.} We don't execute the graph, so we can't yet
  report actual end-to-end latency, cost, or answer accuracy --- only the
  structural predictors of them.
  \item \textbf{Single judge call per decomposition}, no self-consistency;
  LLM-judge scores have variance.
\end{itemize}

\section{Future work (1 min)}

\say
\begin{itemize}
  \item \textbf{Execute the DAG} --- answer sub-queries with the routed models,
  measure real latency, real cloud cost, and final-answer accuracy against
  multi-hop QA datasets (HotpotQA, MuSiQue). Then correlate our structural
  predictors (depth, width) with measured performance.
  \item \textbf{Judge ensemble} to kill self-bias.
  \item \textbf{Add the trained checkpoints} for R2-Reasoner / HybridFlow to close
  the reproduction gap.
  \item \textbf{A recommender} --- given a new query, predict which decomposer to
  use.
  \item \textbf{Optional gold set} --- hand-annotate $\sim$50 queries so we can
  report recall/precision against references alongside the reference-free scores.
\end{itemize}

\section{Anticipated questions + answers}

\paragraph{Q: What's novel here --- you didn't invent a decomposition algorithm.}
The novelty is the evaluation. Three things: (1)~a normalization layer that maps
heterogeneous method outputs --- XML, JSON, prose --- onto one comparable
sub-query\,+\,DAG schema with guaranteed validity; (2)~a reference-free scoring
rubric purpose-built for decomposition quality, since no gold-standard exists;
(3)~the controlled baseline design, especially the same-prompt/different-scale
pair that separates method contribution from model size. It's a
systems-and-evaluation contribution, not an algorithms one.

\paragraph{Q: Why not just use a gold-standard dataset and compute exact-match?}
Because ``the correct decomposition'' isn't well-defined for open-ended queries
--- annotators disagree. Datasets like Break/QDMR exist for narrow QA but don't
generalize. Reference-free judging lets us evaluate any query; the cost is that we
measure soundness, not correctness. Adding a small gold set is future work.

\paragraph{Q: How do you know the LLM judge is reliable?}
We don't fully --- that's a stated limitation. We reduce the risk by (a)~computing
everything objective --- counts, depth, width, acyclicity --- deterministically
instead of asking the judge, (b)~giving the judge a tight 1--5 rubric with
explicit criteria rather than a vibe score, (c)~temperature~0. Self-consistency
and a judge ensemble are the next step.

\paragraph{Q: Why does the judge share the large model with some decomposers?}
Practicality --- one model to host. It biases toward large-tier methods; we flag
it in the UI and the judge model is a config setting. Proper fix is a judge from a
different model family.

\paragraph{Q: The scoreboard resets on reload --- why no database?}
Deliberate. The tool is for interactive exploration, not a persistent benchmark
leaderboard. Statelessness keeps it simple and reproducible. A persistent
benchmark mode is future work once we add DAG execution.

\paragraph{Q: What does ``parallelism / width'' actually mean formally?}
The size of the maximum antichain in the dependency DAG --- the largest set of
sub-queries with no path between any two of them. By Dilworth's theorem that
equals the minimum number of chains needed to cover the graph, which is the
minimum number of sequential ``waves'' of execution. It's the theoretical max
speedup from parallel execution.

\paragraph{Q: Could a method game your metrics?}
The structural ones, somewhat --- you could emit width-heavy graphs to look
parallelizable. That's exactly why dependency-correctness is a judged criterion:
the judge checks whether the claimed parallel edges are actually justified.
Structural metrics describe the graph; the judge checks whether the graph is
\emph{honest}.

\vfill
\hrulefill

\noindent\footnotesize
Sources: HybridFlow --- arXiv:2512.22137 (ICML 2026);
Route-and-Reason / R2-Reasoner --- arXiv:2506.05901 (WWW 2026);
Uno-Orchestra --- arXiv:2605.05007;
Least-to-Most --- arXiv:2205.10625 (Zhou et al., 2022);
Self-Ask --- arXiv:2210.03350 (Press et al., 2022).

\end{document}
